\section{Arquitectura del modelo: MoE vs Dense}

\subsection{Ventajas a corto plazo y riesgos a largo plazo de MoE}

El moderador planteó: MoE y la atención dispersa ya son prácticas dominantes; ¿la atención del grupo a la eficiencia ya superó a la atención a los resultados?

\begin{importantbox}{Xue Fuzhao (薛福昭): MoE podría ser eliminado a largo plazo}
En esencia, MoE intercambia memory (la cantidad de parámetros) por FLOPs, lo cual a su vez implica una data efficiency menor. Pero \textbf{a largo plazo los FLOPs siempre son lo más barato}: aunque la ley de Moore ya no es perfecta, la velocidad de mejora de los chips sigue siendo mucho más rápida que el crecimiento de los datos. Por lo tanto:
\begin{itemize}
  \item Todo esquema que intercambia otra cosa por FLOPs es una \textbf{solution a corto plazo}
  \item Todo esquema que intercambia FLOPs por otra cosa es una \textbf{solution a largo plazo}
\end{itemize}
El modelo podría regresar gradualmente a Dense, e incluso convertirse en Super-Dense, como el Diffusion Model.
\end{importantbox}

\textbf{Cao Yu (曹玉) (perspectiva de posentrenamiento)}: Está completamente de acuerdo en que los FLOPs son la única variable determinista. Pero añade un punto: el preentrenamiento tradicional sin datos sintéticos «posiblemente ya esté dead». Modelos como Gemini incorporan una gran cantidad de datos sintéticos desde la etapa de preentrenamiento, lo cual equivale a \textbf{intercambiar FLOPs por Data}, y luego usar Data para hacer Scaling. Si el preentrenamiento mismo ya está haciendo «FLOPs for Data», entonces MoE, como \textbf{solución de ingeniería dentro de un sistema de hardware limitado por el ancho de banda}, todavía tiene su valor.

\textbf{Liu Ziwei (刘子伟)}: Existe un Trade-off entre el corto y el largo plazo; personalmente se inclina un poco más hacia Dense Model. Pero, considerando escenarios concretos de despliegue, MoE todavía tiene sus características de uso propias. Además, hay quienes exploran la próxima generación de arquitecturas de cómputo (con un acoplamiento más estrecho entre almacenamiento y cómputo), lo cual podría influir en la elección de arquitectura dentro de 3 a 5 años.

\textbf{Xie Tianbao (谢天宝)}: La investigación en Model Architecture ya entró en una etapa de gran floración de propuestas. MoE es solo el esquema maduro de la industria; en el ámbito académico todavía hay mucha más exploración. Lo importante es que \textbf{el diseño de la arquitectura debe seguir al Learning Paradigm}: si en el futuro el Continual Learning se vuelve importante, la elección de arquitectura también cambiará en consecuencia.

\begin{warningbox}{El consenso de OpenAI: de Compute-Bound a Data-Bound}
Xue Fuzhao señaló que OpenAI declaró públicamente, en el lanzamiento de GPT-4.5, que en general ya se había pasado de Compute-Bound a Data-Bound: este es el consenso de toda la Community.
\end{warningbox}

\subsection{Resumen de la sección}
MoE es la solución de ingeniería óptima bajo las restricciones actuales de hardware, pero a largo plazo podría regresar a Dense. El paradigma de preentrenamiento ya cambió: la incorporación masiva de datos sintéticos ha difuminado los límites del «preentrenamiento tradicional». La elección de arquitectura depende, en última instancia, de la evolución del hardware y del paradigma de aprendizaje.

%% ===== Multimodalidad nativa =====
